\begin{helper}
Let \(L=\log n\), \(q=n^{1/4}\), and
\[
k_{\mathbb R}=(1+\eta)\sqrt{nL}.
\]
Assume \(\eta>0\), \(\varepsilon_1>0\), \(n>0\), \(L>0\), \(k\ge0\),
\[
k_{\mathbb R}\le k,\qquad
\frac{2}{L}\le\frac{\varepsilon_1}{2},
\]
and
\[
\frac{1}{(1+\eta)q\sqrt L}\le\frac{\varepsilon_1}{2}.
\]
Then
\[
k\left(\frac{2k}{L}+q\right)\le \varepsilon_1 k^2.
\]
\end{helper}
\begin{proof}
The first threshold gives
\[
k\frac{2k}{L}=\frac{2}{L}k^2
  \le\frac{\varepsilon_1}{2}k^2.
\]
For the second term, multiply
\[
\frac{1}{(1+\eta)q\sqrt L}\le\frac{\varepsilon_1}{2}
\]
by the positive denominator to get
\[
1\le\frac{\varepsilon_1}{2}(1+\eta)q\sqrt L.
\]
Multiplying by \(q\ge0\) gives
\[
q\le \frac{\varepsilon_1}{2}(1+\eta)q^2\sqrt L.
\]
Since \(q^2=\sqrt n\), the right hand side is
\[
\frac{\varepsilon_1}{2}(1+\eta)\sqrt n\sqrt L
  =\frac{\varepsilon_1}{2}k_{\mathbb R}.
\]
Using \(k_{\mathbb R}\le k\), we obtain
\[
q\le\frac{\varepsilon_1}{2}k.
\]
Multiplying by \(k\ge0\) yields
\[
kq\le\frac{\varepsilon_1}{2}k^2.
\]
Adding the two half-bounds proves
\[
k\left(\frac{2k}{L}+q\right)
=k\frac{2k}{L}+kq
\le\frac{\varepsilon_1}{2}k^2+\frac{\varepsilon_1}{2}k^2
=\varepsilon_1 k^2.
\]
\end{proof}
